\section{Related work}
\label{sec:related}

\subsection{Calibration and fusion with correlated measurements}

Field calibration learns a relationship between sensor responses and a reference quantity under the conditions in which the instruments are used~\citep{devito2008airquality,zimmerman2018calibration}. Drift introduces an additional challenge because that relationship can change over time~\citep{vergara2012drift}. Our setting uses a short labeled calibration record for each deployment and asks for predictions after the reference target is hidden. It does not estimate a universal response curve for a named instrument. The support must instead identify the response and dependence structure needed for the current episode.

Correlated errors have also been studied in distributed estimation. Covariance intersection addresses fusion when correlations between estimates are unknown~\citep{julier1997ci}. More recent neural approaches include IFNet and DIFNet for fusion with unknown correlations in measurement noise~\citep{wang2025ifnet,dong2025difnet}. DIFNet also considers changes in sensor nodes and network topology. KalmanNet combines learned components with filtering for partially known dynamics~\citep{revach2022kalmannet}. These approaches establish the importance of dependence aware fusion. The distinction here is the episodic calibration problem: sensor responses are inferred from labeled support rows, and the query is a contemporaneous subset of readings rather than a step in a specified state transition model.

\subsection{Products of experts and learned approximate inference}

A product of experts combines factors by multiplication~\citep{hinton2002poe}. In multimodal models, omitting a factor offers a convenient way to handle an absent modality~\citep{wu2018mvae,ma2021smil}. Multivariate latent products are already part of that literature, so lifting the prediction variable should not be read as introducing the general idea of a latent product. Our analysis concerns a narrower question: which variables must remain in a static Gaussian site representation if it is to reproduce sensor conditionals under shared noise? The scalar obstruction and the dimension result identify a limitation of target only factors under explicit assumptions. They do not exclude scalar sites that are recomputed from the full observation context.

Expectation propagation refines an approximate factor using a cavity distribution formed by removing the current factor~\citep{minka2001ep}. Learned message operators can amortize those updates~\citep{heess2013learning,eslami2014just,jitkrittum2015kernel}. Our small model uses this organization: it removes one sensor's site, summarizes what the other sensors imply about target and nuisance, and predicts a replacement site. The unknown sensor response is inferred from calibration examples rather than supplied as a known likelihood. As a result, the model is inspired by expectation propagation but is trained by predictive likelihood, without an explicit moment matching constraint at every update.

Posterior linearization provides another connection. A nonlinear response can be approximated by a local linear Gaussian measurement under a current posterior approximation~\citep{deisenroth2012ep,garciafernandez2015plf}. We use the same interpretable quantities, namely a response offset, a target slope, nuisance loadings, and noise variance. The network learns those quantities from support and cavity statistics. This also follows the broader idea of algorithm unrolling, in which a finite sequence of inference operations is combined with trainable updates~\citep{gregor2010lista,monga2021unrolling}. The resulting guarantees are structural, such as positive posterior variance and insensitivity to hidden readings, rather than convergence guarantees for an inference algorithm.

\subsection{Prediction with incomplete inputs}

The observation process determines whether missingness can be ignored in inference~\citep{rubin1976inference,little2019statistical}. Supervised prediction with missing inputs also has statistical requirements that depend on the relationship between masks, inputs, and targets~\citep{josse2019consistency}. These considerations remain relevant even when a model is defined for every mask. A prediction rule that omits absent likelihood factors is correct under its measurement model only when the mask supplies no further information needed for that conditional. Informative selection requires a model for selection or an explicitly different prediction objective.

NeuMiss derives a neural architecture from the Gaussian predictor with incomplete inputs~\citep{lemorvan2020neumiss}. It shares our interest in letting a statistical calculation guide network design. Its usual setting trains a predictor for a dataset with specified feature identities. Our model is trained across tasks and obtains each sensor's calibration from the current support. This distinction motivates shared sensor operations and a readout whose matrix dimension depends on the latent nuisance dimension rather than the number of available sensor columns.

\subsection{Learning from a support set}

Conditional and attentive neural processes learn predictive distributions conditioned on sets of examples~\citep{garnelo2018cnp,kim2019anp}. DeepSets and attention over sets provide reusable ways to encode unordered collections~\citep{zaheer2017deepsets,lee2019settransformer}. The compact model uses a DeepSets encoder for the observed calibration pairs of one sensor. Its context exchange then occurs through Gaussian cavity statistics, which preserve an explicit target and nuisance interpretation.

Prior fitted networks train on tasks drawn from a prior so that a forward pass approximates inference for a new dataset~\citep{muller2022pfn}. TabPFN and its later tabular foundation model extend this approach to tabular prediction~\citep{hollmann2023tabpfn,hollmann2025tabpfnv2}; the cited TabICL work studies classification~\citep{qu2025tabicl}. Our transformer baseline adopts a cell token design inspired by TabPFN v2 and is trained on the sensor task distribution used in this study. It is not an evaluation of the published pretrained TabPFN model. The comparison between its direct Gaussian head and the lifted site head asks whether an explicit fusion calculation helps when the context encoder and training distribution are held fixed. Because the two heads also use different initializations, that comparison measures the complete readout choice rather than isolating a single source of improvement.
